% !TITLE 送杜少府之任蜀州
% !AUTHOR 王勃
% !DYNASTY 唐
% !GENRE 五言律诗
% !ORDER 3

\begin{poem}
城阙\textsuperscript{1}辅三秦\textsuperscript{2}，风烟望五津\textsuperscript{3}。
与君离别意，同是宦游\textsuperscript{4}人。
海内存知己，天涯若比邻\textsuperscript{5}。
无为在歧路\textsuperscript{6}，儿女共沾巾\textsuperscript{7}。
\end{poem}

\begin{annotations}
\notetext{1}{城阙：城门两边的楼观，代指京城长安。阙（què），宫门前的高台建筑。}
\notetext{2}{三秦：指关中地区（今陕西一带）。项羽灭秦后，将秦地分为三部分封给三个秦降将，故称“三秦”。这里泛指长安附近的关中大地。}
\notetext{3}{五津：岷江的五个渡口，在今四川境内。代指蜀州（杜少府即将赴任之地）。}
\notetext{4}{宦游：为了做官而远游他乡。王勃自己也因为做官而长期漂泊，所以说“同是宦游人”。}
\notetext{5}{比邻：近邻。只要天下有知心朋友，即使相隔天涯也像近邻一样。这句后来成为中国人表达友情的最高格言。}
\notetext{6}{歧路：岔路口，指分手的地方。}
\notetext{7}{沾巾：泪水沾湿了手巾。不要在岔路口像小儿女一样哭哭啼啼——这是劝慰，也是自勉。}
\end{annotations}

\begin{appreciation}
《送杜少府之任蜀州》是初唐最著名的送别诗之一。王勃写这首诗时还不到二十岁，送的是姓杜的朋友去蜀州（今四川一带）做官。一个不到二十岁的年轻人，劝另一个年轻人别哭，还劝得这么有底气，这在中国送别诗里是头一回。

“城阙辅三秦，风烟望五津。”城阙指京城长安；辅是护卫的意思，长安被三秦大地簇拥着。五津是岷江的五个渡口，代指朋友要去的地方。人还没走，先望向远方：长安和蜀州隔着万水千山，望过去只有风烟茫茫。“与君离别意，同是宦游人”——宦游，就是为做官而远走他乡。你我都是离家在外的人，离别的滋味，不用多说。

“海内存知己，天涯若比邻。”还记得我们读过的《行行重行行》吗？那里头有个女子，和游子“各在天一涯”，觉得天涯是思念的尽头，想得衣带渐宽。同样是“天涯”这个词，王勃把它重新量了一遍：不是距离的尽头，而是友情的刻度。只要这世上还有懂你的人，天涯海角也像隔壁邻居。从此中国人说友谊，最狠的一句话就是这句。

最后是劝告：“无为在歧路，儿女共沾巾。”歧路是岔路口，古人送别常在岔路口分手；沾巾是眼泪打湿手巾。别在岔路口哭哭啼啼。这是劝朋友，也是劝自己——这一劝，就把整首诗提上去了。

王勃后来只活了二十六岁，溺水而亡。但这首诗把少年人那种“我信友情能战胜距离”的底气留了下来。你以后和朋友各奔东西，分到不同的班级、不同的城市，如果还能想起“天涯若比邻”这五个字，就说明你交到了值得交的朋友。
\end{appreciation}
